\scanpage{121}
\[
\deleted{\pi(x)=\varphi(a,x),}\tag*{\deleted{(11)}}
\]
Ce qu'on obtient en faisant $x=(a)$ dans (9), et notant que $\varphi(a,a)=2q(a)=2$

\inserted{\uline{Scholie} La cat des $(P,Q,H)$ comme dessus équivaut donc à celle des $(T,q,a)$, $T$ fibré vectoriel de \unsure{rg|rang} $2\nu$, $q$ forme quadratique \unsure{lisse} sur $T$, $a$ section de $T$, tels que $q(a)=1$}

\uline{NB} a) En car. résiduelles $\ne2$, on voit donc que $T\simeq V+D\simeq V+\mathcal O_S$ canoniquement, \inserted{(i.e. $E$ est muni d'une section origine canonique)} la décomposition étant ortho et $Q$ étant défini canoniquement par une forme
\[
q(x,t)=q_0(\unsure{x})+t^2\tag*{(11)}
\]
où $q_0$ est une forme quadratique sur $V$ ; $Q\cap E$ est défini par l'équation (en $x\in\Gamma(,E)$)
\[
q_0(x)+1=0\qquad\text{i.e.}\quad q_0(x)=-1\tag*{(12)}
\]
Dans les car. rés $\ne2$, la cat. des triples $(P,Q,H)$, ou des couples $(E,Q\cap E)$ ($E$ étant muni de sa structure de fibré affine \ldots) est \inserted{donc} équivalente à celle des couples $(V,q_0)$, où $V$ est un fibré vectoriel de rang $2\nu-1$.

En car. 2, par contre, on doit avoir automatiquement $D\subset V$, et il n'y a pas\deleted{\unsure{,}} de réduction au cas \unsure{d'une|$+$ une} forme de degré $2\nu-1$.

b) Si $T$ est de rang impair $2\nu+1$, donc $V$ de rang pair $2\nu$, les conditions que $Q$ soit
% archive sheet number: 120
